\documentclass[12pt,a4paper]{article}

% Required packages
\usepackage[utf8]{inputenc}
\usepackage[T2A]{fontenc}
\usepackage[russian]{babel}
\usepackage{amsmath, amssymb, amsfonts}
\usepackage{geometry}
\usepackage{xcolor}
\usepackage{booktabs}
\usepackage{tikz}
\usepackage{pgfplots}
\pgfplotsset{compat=1.18}
\usepackage[most]{tcolorbox}

\geometry{left=25mm, right=25mm, top=25mm, bottom=25mm}

% Custom theorem environments using tcolorbox
\tcbset{
    defstyle/.style={fonttitle=\bfseries, colback=blue!5, colframe=blue!70!black},
    theostyle/.style={fonttitle=\bfseries, colback=green!5, colframe=green!50!black},
    remstyle/.style={fonttitle=\bfseries, colback=gray!5, colframe=gray!50!black}
}

\newtcolorbox{definition}[1][]{defstyle, title=Определение #1}
\newtcolorbox{theorem}[1][]{theostyle, title=Теорема #1}
\newtcolorbox{remark}[1][]{remstyle, title=Замечание #1}

\title{\textbf{Теория численных методов для решения ОДУ}\\ \Large Руководство по анализу начально-краевых задач}
\author{Профессор вычислительной математики}
\date{\today}

\begin{document}

\maketitle
\tableofcontents
\newpage

\section{Введение и базовые понятия}

\subsection{Задача Коши (IVP)}
Рассматривается задача Коши (Initial Value Problem, IVP) для системы обыкновенных дифференциальных уравнений (ОДУ) первого порядка:
\begin{equation}
    \label{eq:ivp}
    \vec{y}'(t) = \vec{f}(t, \vec{y}(t)), \quad \vec{y}(t_0) = \vec{y}_0,
\end{equation}
где $t \in [t_0, T]$, $\vec{y}(t) \in \mathbb{R}^d$, а функция $\vec{f}: \mathbb{R} \times \mathbb{R}^d \to \mathbb{R}^d$ предполагается достаточно гладкой и удовлетворяющей условию Липшица по переменной $\vec{y}$ для гарантии существования и единственности решения (Теорема Пикара — Линделёфа).

\subsection{Теория аппроксимации и ошибки}
Для численного решения введем сетку $t_n = t_0 + n h$, где $h$ --- шаг интегрирования. Обозначим точное решение как $\vec{y}(t_n)$, а его численное приближение как $\vec{y}_n$.

\begin{definition}[Локальная и глобальная погрешности]
\textbf{Локальная погрешность усечения (Local Truncation Error, LTE)} $\vec{d}_{n+1}$ — это ошибка, возникающая за один шаг метода в предположении, что предыдущие значения были точными ($\vec{y}_n = \vec{y}(t_n)$).
\\
\textbf{Глобальная погрешность (Global Error)} $\vec{e}_n = \vec{y}_n - \vec{y}(t_n)$ — это накопленная ошибка к моменту времени $t_n$.
\end{definition}

Метод имеет \textbf{порядок аппроксимации $p$}, если для достаточно гладкой функции $\vec{f}$ локальная погрешность удовлетворяет условию:
\[
    \|\vec{d}_{n+1}\| = \mathcal{O}(h^{p+1}) \quad \text{при } h \to 0.
\]
(Деление на шаг $h$ в некоторых определениях невязки дает $\mathcal{O}(h^p)$).

\subsection{Сходимость и устойчивость}
\begin{theorem}[Теорема эквивалентности (аналог теоремы Лакса/Далквиста)]
Сходимость численного метода ($\vec{e}_n \to 0$ при $h \to 0$) имеет место тогда и только тогда, когда метод \textbf{аппроксимирует} задачу (согласованность, $p \ge 1$) и является \textbf{нуль-устойчивым}. Устойчивость ограничивает рост возмущений, возникающих на каждом шаге.
\end{theorem}

\section{Одношаговые методы (Семейство Рунге-Кутты)}

\subsection{Общий вид метода}
Методы Рунге-Кутты (РК) с $s$ стадиями вычисляют значение на новом слое с помощью промежуточных оценок вектора производных:
\begin{align}
    \vec{k}_i &= \vec{f}\left( t_n + c_i h, \, \vec{y}_n + h \sum_{j=1}^s a_{ij} \vec{k}_j \right), \quad i = 1, \dots, s, \\
    \vec{y}_{n+1} &= \vec{y}_n + h \sum_{i=1}^s b_i \vec{k}_i.
\end{align}

\subsection{Таблицы Бутчера}
Коэффициенты метода удобно записывать в виде \textbf{таблицы Бутчера}:
\[
\renewcommand{\arraystretch}{1.2}
\begin{array}{c|c}
    \vec{c} & A \\
    \hline
            & \vec{b}^T
\end{array}
\quad = \quad
\begin{array}{c|ccc}
    c_1 & a_{11} & \dots & a_{1s} \\
    \vdots & \vdots & \ddots & \vdots \\
    c_s & a_{s1} & \dots & a_{ss} \\
    \hline
        & b_1 & \dots & b_s
\end{array}
\]
Здесь $A \in \mathbb{R}^{s \times s}$, $\vec{b}, \vec{c} \in \mathbb{R}^s$. Обычно выполняется условие согласования узлов: $c_i = \sum_{j=1}^s a_{ij}$.

\subsection{Условия порядка}
Для достижения порядка $p$ коэффициенты метода должны удовлетворять системе алгебраических уравнений (условия деревьев). Для методов малых порядков они имеют вид:
\begin{itemize}
    \item \textbf{Порядок 1 ($p=1$):} $\sum b_i = 1$ (1 уравнение).
    \item \textbf{Порядок 2 ($p=2$):} Добавляется $\sum b_i c_i = \frac{1}{2}$ (всего 2).
    \item \textbf{Порядок 3 ($p=3$):} Добавляются $\sum b_i c_i^2 = \frac{1}{3}$ и $\sum b_i a_{ij} c_j = \frac{1}{6}$ (всего 4).
    \item \textbf{Порядок 4 ($p=4$):} Добавляются еще 4 уравнения (всего 8), в том числе:
    \[
        \sum b_i c_i^3 = \frac{1}{4}, \quad \sum b_i c_i a_{ij} c_j = \frac{1}{8}, \quad \sum b_i a_{ij} c_j^2 = \frac{1}{12}, \quad \sum b_i a_{ij} a_{jk} c_k = \frac{1}{24}.
    \]
\end{itemize}

\subsection{Классификация методов}
Структура матрицы $A$ определяет тип метода:
\begin{itemize}
    \item \textbf{Явные (ERK):} $a_{ij} = 0$ при $j \ge i$. Матрица $A$ строго нижнетреугольная. Стадии $\vec{k}_i$ вычисляются последовательно.
    \item \textbf{Неявные (IRK):} Матрица $A$ заполнена. Требуется решение системы нелинейных алгебраических уравнений размерности $s \times d$.
    \item \textbf{Диагонально неявные (DIRK):} $a_{ij} = 0$ при $j > i$. Матрица нижнетреугольная. Стадии вычисляются последовательно, но на каждой стадии решается система размерности $d$.
\end{itemize}

\section{Теория устойчивости (Фундамент)}

\subsection{Модельное уравнение Далквиста}
Линейный анализ устойчивости базируется на \textbf{тестовом уравнении Далквиста}:
\begin{equation}
    y' = \lambda y, \quad \lambda \in \mathbb{C}, \quad \text{Re}(\lambda) < 0.
\end{equation}

\begin{remark}[Обоснование применимости]
Для произвольной нелинейной системы $\vec{y}' = \vec{f}(\vec{y})$ поведение решений вблизи положения равновесия $\vec{y}^*$ определяется линеаризацией $\vec{y}' \approx J (\vec{y} - \vec{y}^*)$, где $J = \frac{\partial \vec{f}}{\partial \vec{y}}$. Если матрица $J$ диагонализуема ($J = V \Lambda V^{-1}$), система распадается на $d$ независимых скалярных уравнений вида $u_i' = \lambda_i u_i$, где $\lambda_i$ — собственные значения якобиана. Таким образом, анализ одного скалярного комплексного уравнения охватывает динамику всей линейной системы.
\end{remark}

\subsection{Устойчивость одношаговых методов}
Применим общий метод РК к уравнению $y' = \lambda y$. Получаем систему:
$k_i = \lambda \left( y_n + h \sum a_{ij} k_j \right)$ и $y_{n+1} = y_n + h \sum b_i k_i$.
Вводя вектор $\vec{k} = (k_1, \dots, k_s)^T$, шаг $z = \lambda h$, и $\vec{1} = (1, \dots, 1)^T$, запишем:
\[
    \vec{k} = z y_n \vec{1} + z A \vec{k} \implies \vec{k} = z (I - zA)^{-1} \vec{1} y_n.
\]
Подставляя в формулу шага, получаем $y_{n+1} = R(z) y_n$, где:
\begin{equation}
    R(z) = 1 + z \vec{b}^T (I - zA)^{-1} \vec{1}.
\end{equation}
Функция $R(z)$ называется \textbf{функцией устойчивости}. Для ERK $R(z)$ является полиномом степени не выше $s$. Для IRK это рациональная функция.

\begin{definition}[Область устойчивости]
Область абсолютной устойчивости $S$ определяется как множество значений $z$, при которых численное решение не нарастает:
\[
    S = \{ z \in \mathbb{C} \mid |R(z)| \le 1 \}.
\]
\end{definition}

% ТikZ Placeholder для Области Устойчивости
\begin{figure}[h]
    \centering
    \begin{tikzpicture}
        \begin{axis}[
            width=8cm, height=8cm,
            axis lines=middle,
            xmin=-3, xmax=1, ymin=-2, ymax=2,
            xlabel=$\text{Re}(z)$, ylabel=$\text{Im}(z)$,
            title={Пример области устойчивости $S$ (Явный метод Эйлера)},
            xtick={-2,-1,0}, ytick={-1,1},
            grid=both
        ]
            % Placeholder for a stability region (Circle of radius 1 at x=-1)
            \draw[fill=blue!20, opacity=0.5, thick] (axis cs:-1,0) circle [radius=1];
        \end{axis}
    \end{tikzpicture}
    \caption{Область абсолютной устойчивости явного метода Эйлера (единичный круг $|1+z| \le 1$). Для неявных методов область $S$ часто охватывает левую полуплоскость.}
\end{figure}

\subsection{Устойчивость многошаговых методов}
Многошаговые методы связывают несколько слоев:
\begin{equation}
    \sum_{j=0}^k \alpha_j y_{n+j} = h \sum_{j=0}^k \beta_j f_{n+j}.
\end{equation}
Для анализа $y'=\lambda y$ введем характеристические полиномы $\rho(\xi) = \sum \alpha_j \xi^j$ и $\sigma(\xi) = \sum \beta_j \xi^j$. Подстановка $y_n = \xi^n$ дает характеристическое уравнение:
\begin{equation}
    \pi(\xi, z) \equiv \rho(\xi) - z \sigma(\xi) = 0, \quad z = \lambda h.
\end{equation}

\begin{definition}[Корневое условие Далквиста]
Для нуль-устойчивости (при $h \to 0$, то есть $z=0$) корни полинома $\rho(\xi)$ должны удовлетворять условию: $|\xi_i| \le 1$, причем корни на единичной окружности ($|\xi_i| = 1$) должны быть простыми.
\end{definition}

\subsection{Жесткость (Stiffness) и $A$-устойчивость}
\begin{definition}[Жесткая система]
Система называется жесткой, если ее якобиан имеет собственные значения с сильно различающимися отрицательными вещественными частями (число обусловленности велико), при этом быстрые затухающие компоненты диктуют малый шаг для явных методов из-за условий устойчивости (число Куранта/CFL-подобные ограничения).
\end{definition}

\begin{definition}
Метод является \textbf{$A$-устойчивым}, если его область устойчивости $S$ содержит всю левую комплексную полуплоскость: $\mathbb{C}^- \subseteq S$. 
Метод называется \textbf{$L$-устойчивым}, если он $A$-устойчив и $\lim_{|z|\to\infty} R(z) = 0$.
\end{definition}
Ни один явный метод РК или многошаговый метод не может быть $A$-устойчивым. Для жестких задач необходимы неявные методы.

\section{Многошаговые и специальные методы для жестких систем}

\subsection{Методы Адамса}
Методы Адамса строятся интегрированием $\vec{y}'=\vec{f}$ и заменой $\vec{f}$ интерполяционным полиномом.
\begin{itemize}
    \item \textbf{Адамс-Башфорт (явные):} $\beta_k = 0$. Пример 2-го порядка: 
    $y_{n+2} = y_{n+1} + \frac{h}{2}(3f_{n+1} - f_n)$.
    \item \textbf{Адамс-Моултон (неявные):} $\beta_k \ne 0$. Обладают большей областью устойчивости. Пример (правило трапеций): 
    $y_{n+1} = y_n + \frac{h}{2}(f_{n+1} + f_n)$.
\end{itemize}

\subsection{Методы Гира (BDF - Формулы обратного дифференцирования)}
Вместо аппроксимации функции $f$, методы BDF аппроксимируют само решение $y(t)$ интерполяционным полиномом по узлам $t_{n+k}, t_{n+k-1}, \dots$ и дифференцируют его в точке $t_{n+k}$.
\begin{equation}
    \sum_{j=0}^k \alpha_j y_{n+j} = h \beta_k f_{n+k}.
\end{equation}
\textbf{Свойства BDF:} Они являются де-факто стандартом для жестких систем. Метод BDF1 — это неявный метод Эйлера (L-устойчив). BDF2 также L-устойчив. При $k=3,4,5,6$ методы $A(\alpha)$-устойчивы, а для $k \ge 7$ они становятся нуль-неустойчивыми (первый барьер Далквиста).

\subsection{Методы Розенброка (W-методы)}
Неявные методы требуют решения нелинейной системы на каждом шаге (обычно методом Ньютона). \textbf{Методы Розенброка} избегают этого, встраивая одну итерацию метода Ньютона прямо в формулы стадий.
Общий вид метода Розенброка:
\begin{align}
    (I - \gamma h J) \vec{k}_i &= h \vec{f}\left( \vec{y}_n + \sum_{j=1}^{i-1} a_{ij} \vec{k}_j \right) + h J \sum_{j=1}^{i-1} \gamma_{ij} \vec{k}_j, \\
    \vec{y}_{n+1} &= \vec{y}_n + \sum_{i=1}^s b_i \vec{k}_i,
\end{align}
где $J = \frac{\partial \vec{f}}{\partial \vec{y}}(\vec{y}_n)$ — якобиан, вычисленный один раз в начале шага. На каждой стадии решается \textbf{линейная} система с одной и той же матрицей $(I - \gamma h J)$, что требует лишь одного LU-разложения.

\section{Продвинутая тема: Метод Гира в представлении Нордсика}

Для реализации методов с переменным шагом и переменным порядком (что критически важно в современных решателях типа CVODE/LSODE) стандартная форма BDF неудобна: при изменении $h$ коэффициенты $\alpha_j, \beta_k$ сложным образом пересчитываются. Нордсик (Nordsieck, 1962) предложил эквивалентную одношаговую формулировку.

\subsection{Вектор Нордсика}
Вместо хранения истории значений $y_n, y_{n-1}, \dots, y_{n-k+1}$, хранится аппроксимация решения и его производных, масштабированных шагом $h$:
\begin{equation}
    \vec{z}_n = \begin{pmatrix} y_n \\ h y'_n \\ \frac{h^2}{2!} y''_n \\ \vdots \\ \frac{h^k}{k!} y^{(k)}_n \end{pmatrix} \in \mathbb{R}^{d \times (k+1)}.
\end{equation}

\subsection{Алгоритм: Предиктор-Корректор}
\textbf{1. Предиктор (прогноз):} Переход на новый слой основывается на разложении Тейлора. В векторной форме это линейное преобразование:
\[
    \vec{z}_{n(0)} = P \vec{z}_{n-1},
\]
где $P$ — верхнетреугольная матрица Паскаля. Например, для $k=2$:
\[
    P = \begin{pmatrix} 1 & 1 & 1 \\ 0 & 1 & 2 \\ 0 & 0 & 1 \end{pmatrix}.
\]

\textbf{2. Корректор (уточнение):} Оценивается невязка прогноза:
\[
    \vec{e}_n = h \vec{f}(t_n, \vec{y}_{n(0)}) - h \vec{y}'_{n(0)}.
\]
Затем вектор Нордсика обновляется с помощью вектора коэффициентов метода $\vec{l}$:
\[
    \vec{z}_n = \vec{z}_{n(0)} + \vec{l} \otimes \vec{e}_n.
\]
Вектор $\vec{l}$ однозначно связан с коэффициентами многошагового метода $\alpha_j, \beta_k$.

\subsection{Преимущества для переменного шага}
Если необходимо изменить шаг интегрирования с $h_{old}$ на $h_{new}$, вводится отношение $r = h_{new}/h_{old}$. Вектор Нордсика пересчитывается простым диагональным масштабированием без потери точности и без сложной полиномиальной интерполяции сетки:
\begin{equation}
    \vec{z}_{n} \leftarrow C(r) \vec{z}_n, \quad \text{где } C(r) = \begin{pmatrix} 1 & 0 & 0 & \dots & 0 \\ 0 & r & 0 & \dots & 0 \\ 0 & 0 & r^2 & \dots & 0 \\ \vdots & \vdots & \vdots & \ddots & \vdots \\ 0 & 0 & 0 & \dots & r^k \end{pmatrix}.
\end{equation}
Такое элегантное свойство делает методы BDF в форме Нордсика основой подавляющего большинства промышленных программных пакетов для жестких ОДУ.

\end{document}